\label{cap:introducao}

Exames de diagnóstico por imagem, como ultrassonografia e radiologia, são procedimentos comuns no cotidiano popular e na rotina de saúde, como evidenciado na notícia de \citeonline{MedicinaSA2024}, que aponta um total de 2,4 bilhões de exames diagnósticos realizados. O setor clínico de diagnósticos vêm se tornando um alicerce não apenas da saúde, como tendência médica em análises de caso, mas também um pilar econômico, com geração de receita bruta de R\$ 24 bilhões movimentados no ano de 2022 \cite{Abramed_2023}, sendo metade desse valor realizado pela saúde privada.
Autores como \citeonline{Yasui_Fujiwara_Okitsu_Yonemitsu_Ito_Yokosuka_2012} e \citeonline{OBRIEN_2007} avaliam a importância do uso de diagnóstico por exames para constatação precoce de doenças e elevação da qualidade no acompanhamento médico em detrimento ao custo relativo de tais operações, ponto comum de discussão, uma vez que os custos operacionais para uso e manutenção da aparelhagem necessária é elevado. Entretanto, como aponta \citeonline{CBDL_2024}, o custo da operação vêm caindo anualmente, como exemplificada pela queda de 22,9\% para 21\% no ano de 2023, influindo diretamente nos preços dos exames e estabilidade funcional.

Apesar de sua importância, profissionais capacitados nas áreas de radiologia e ultrassonografia estão em falta para operação da aparelhagem. Com uma estimativa de crescimento de 6\% para 18\% nas vagas não preenchidas por especialistas na área, espera-se um grande crescimento nas filas de atendimento dessa modalidade, devido o envelhecimento populacional global e as decorrentes necessidades de acompanhamento na saúde\cite{RSNA_2024}. No Brasil, verifica-se uma realidade de escassez generalizada para profissionais da área: para cada 1,99 milhões de indivíduos estimam-se apenas 1,51 mil médicos no ano de 2017, como apontam \citeonline{Liebel_Dias_Schneider_SaScielo2021}.

O excedente de demanda de exames pelos pacientes e a carência de disponibilidade de atendimento, causada pela escassez de mão de obra, geram custos na operação devido ao maquinário parado e ao afunilamento no atendimento. Além disso, segundo relatório público da \citeonline{ANS_2024}, embora o prejuízo fiscal da saúde suplementar tenha reduzido significativa, as margens operacionais continuam próximas ao limiar do adequado, pressionando as empresas do ramo a procurar formas de cortar custos.

Uma possibilidade para o controle de gastos, através do incremento na qualidade operacional, é o escalonamento da força de trabalho e a elaboração de escalas médicas otimizadas para atender altas demandas por equipes reduzidas. Autores como \citeonline{Rousseau_Pesant_Gendreau_2002} apontam que a geração de quadros de horário com boa qualidade podem aumentar significativamente a capacidade de atendimento de pacientes, promovendo a redução nas filas de espera, aumentando a satisfação do cliente, tudo isso enquanto mantém o tamanho da equipe baixa. Além disso, \citeonline{Robinson_Chen_2003} citam também a importância do processo para melhor aproveitamento do maquinário disponível, reduzindo tempo ocioso do equipamento e, impulsionando a operação. 
Entretanto, embora altamente benéfica, a confecção de escalas e quadros de horário, comumente realizada de maneira manual, é uma tarefa laboriosa.

Diante dos pontos discutidos, somados à crescente demanda populacional por serviços médicos de diagnóstico apontada por \citeonline{Liebel_Dias_Schneider_SaScielo2021}, o processo de modelagem matemática e desenvolvimento de algoritmos computacionais para confecção de escalas se torna evidente e necessário. Esta dissertação de mestrado visa estudar o problema de escalonamento da equipe médica em contextos de consultas pré-agendadas, onde se conhece previamente a demanda a ser atendida.

O referido estudo é baseado em dados reais disponibilizados por uma rede de clínicas de saúde suplementar do ramo de diagnóstico por imagem com especialidade em ultrassonografia localizada em uma grande cidade do estado de São Paulo. Observa-se conjunto de dados uma média de 37.000 horários semanais para consulta distribuídos entre 37 unidades, que, em média, disponibilizam até 4 salas individuais para exames, e 294 médicos especialistas associados disponíveis para alocação para realização dos exames. 

Nesse cenário, os profissionais têm uma disponibilidade de horas para realizar exames em um determinado período do dia, além da possibilidade de ausência programada, gerando a necessidade de realocar outro especialista (possivelmente de outra unidade) para suprir a demanda. Contudo, flexibilizar as trocas constantes de profissionais entre unidades para atender a demanda pode resultar em descontentamento da equipe devido a muitas mudanças ao longo do planejamento. Assim, este trabalho visa a alocação de médicos nas unidades de saúde em que atendem, buscando atender os exames e minimizar essas mudanças entre unidades. Todavia, não é comum na literatura acadêmica explorar e estudar modelagens referentes à minimização de trocas de médicos entre unidades em problemas de escalonamento médico, especialmente ao se considerar a possibilidade de múltiplas unidades compartilharem profissionais entre si. Desta forma, investiga-se o problema no contexto de otimização bi-objetivo, buscando alocar os profissionais de modo a minimizar os exames não atendidos enquanto flexibiliza as trocas entre clínicas. No entanto, também se busca minimizar as trocas entre unidades. Para lidar com o problema, foi proposto um modelo de otimização inteira, além de duas abordagens heurísticas baseadas na formulação matemática (Relax-and-Fix e Fix-and-Optimize) e uma meta-heurística evolucionária multiobjetivo (NSGA-II).

Assim, neste trabalho têm-se como objetivos principais: (1) a modelagem e representação matemática do escalonamento da equipe médica, (2) a confecção dos quadros de horário dos médicos por métodos exatos do cenário mediante demanda e disponibilidade já conhecidas através do método determinístico, (3) o estudo e implementação de técnicas de otimização bi-objetivo que explorem a minimização de remanejamentos e trocas de profissionais entre unidades, e, (4) contribuir para a literatura acadêmica e o conhecimento científico através do estudo de sistemas que abordem a alocação em múltiplas unidades ao mesmo tempo que se reduz a troca de clínicas para o profissional ao longo do horizonte de planejamento.

Espera-se que, abordando tais desafios, a dissertação contribua para o aprimoramento do saber científico e acadêmico, bem como melhorias significativas fornecendo modelos sistêmicos bem estruturados para gestão operacional de instituições médicas, tornando-as mais eficientes, impactando diretamente na eficiência do serviço prestado e na satisfação do médico e paciente. 

Para tal, o referido estudo encontra-se estruturado da seguinte maneira: o presente Capítulo \ref{cap:introducao} discorre o contexto a qual o problema está inserido e apresenta a proposta de intervenção pautada no cenário descrito. No Capítulo \ref{cap:revisao-bibliografica} é realizada a Revisão da Literatura em Pesquisa Operacional Aplicada à Saúde enquanto se exploram lacunas no conhecimento acadêmico e metodológico. O Capítulo \ref{cap:problema} descreve de maneira aprofundada o problema observado mediante, descrevendo a característica e classificação do tema estudado. No Capítulo \ref{cap:mip-heuristicas} descreve o modelo matemático proposto para descrever o assunto observado, bem como expandir as diversas técnicas computacionais empregadas para obtenção de resultados. O Capítulo \ref{cap:resultados} por sua vez, apresenta o conjunto de dados e explica o ambiente computacional utilizado para executar a metodologia proposta, ao mesmo tempo, em que expõe e analisa os resultados observados nos experimentos. Por fim, o \ref{cap:fechamento} tem por finalidade apresentar conclusões extraídas do estudo e sugerir melhorias para com o projeto, e possíveis cenários para estudo no tema.